\section{Background}

The hippocampal-entorhinal circuit supports memory retrieval through interactions among sparse coding and inhibitory competition, with temporal coordination providing the organizing frame. In the dentate gyrus and CA3, inhibitory control contributes to pattern separation and winner-take-all dynamics, allowing similar inputs to activate distinct population states \cite{arxiv210506057,arxiv170507614}. In computational models, context-dependent gating and sparse hippocampal memory architectures reduce interference without requiring erasure of prior representations \cite{arxiv171109876,arxiv220200159}.

At the CA1 output stage, the active memory state is not fully captured by firing rate alone. Several retrieved references treat CA1 as a temporal and predictive structure in which spike timing and co-firing relationships interact with replay and phase organization to carry information about sequence and transition structure \cite{arxiv191211126,arxiv200611975,arxiv230511982}. This distinction is important for memory updating. If retrieval selection depends on which context-value association controls CA1 temporal ensemble structure, then disruption of the mechanism should be visible in timing relationships and cross-correlated activity, not only in mean rate changes.

The medial entorhinal cortex provides a major route through which spatial and contextual information reaches CA1. The temporoammonic pathway supplies direct entorhinal input to distal CA1 dendrites, where local inhibition can regulate gain and delay sensitivity while shaping synchrony \cite{arxiv240910460,arxiv240910431}. In parallel, models of the entorhinal-hippocampal loop describe CA1 as a comparison and readout layer that balances internal hippocampal sequences with current entorhinal context \cite{arxiv250112615,arxiv200611975}.

Medial entorhinal inhibitory populations provide a plausible transformation layer between septal input and CA1 output. Medial entorhinal cortex influences CA1 through both direct temporoammonic input and indirect pathways via CA3, providing multiple routes through which inhibitory timing can shape CA1 ensemble structure. Parvalbumin-expressing interneurons and related inhibitory motifs constrain spike timing and support phase locking while shaping the distribution of neural responses in recurrent circuits \cite{arxiv160903622,arxiv220910634,arxiv240517745}. Septal GABAergic projections can therefore be treated as an upstream timing source, with medial entorhinal inhibitory microcircuits providing the local timing transformation and CA1 temporal ensemble organization serving as the retrieval readout that drives behavior.
